\cotefile{83}
\scanpage{129}\mspage{PU 52}
% archive sheet number: 128
% Two ink diagonals cancel the body as a whole. The separate deleted spans follow its paragraphs and displays.

\deleted{Parmi les cinq \deleted{termes} attracteurs yang qui figurent dans le ``paquet''
rappelé à l'instant, nous avions déjà \deleted{\unsure{choi|cho}} mis en védette deux}

\[
\text{\deleted{s\,\deleted{\unsure{x}}i\,m\,p\,l\,i\,c\,i\,t\,é \quad, \quad o\,r\,d\,r\,e \quad,}}
\]

\deleted{auxquels les autres trois apparaissent comme subordonnés. Puis on s'était
aperçu que les qualités qu'ils représentent sont, à vrai dire, ``comme
l'âme et le corps d'une seule et même qualité''. C'est ``le corps'' et non
l'âme, l'aspect le plus yang de cette qualité commune, que je me suis
alors senti incité à mettre en avant, \deleted{à savoir}``l'ordre'' , pour le voir
aussitôt convoler avec son conjoint yin prédestiné, ``l\,e \quad m\,y\,s\,t\,è\,re\handwritten{''}\unsure{.}}

\deleted{En voyant se consommer ces épousailles, \deleted{\unsure{ja}}je ne me suis pas trop laisser
troubler par \handwritten{l}e fait\deleted{\unsure{,}} que dans le couple \deleted{par lequel je venais d'introduire ``l'ordre'', savoir}}

\[
\text{\deleted{ordre -- chaos}}
\]

\deleted{par lequel je venais d'introduire ``l'ordre'' dans notre fameux ``paquet'',
il figure comme époux du c\,h\,a\,o\,s , et nulle\inserted{\handwritten{ment}} du mystère. Je \deleted{ne me} n'ai
fait \deleted{\unsure{à}}attention à la chose qu'après-coup, tant les épousailles de
l'ordre et du mystère semblaient couler de source ! J'y ai resongé par
la suite -- comment il se faisait que ``le chaos'' semblait avoir disparu
de la réflexion. Il n'y en a trace notamment dans la Fleur cosmique. L'ordre
y figure comme le terme yang central, comme le \deleted{\unsure{s}}``sommet'' le plus
élevé, celui autour duquel semblent se grouper tous les autres, avec
les ``ravins'' ou ``points-vallées'' qui leur correspondent. Et ce maître
sommet s'accouple avec ``le mystère'' sur le versant yin, avec ``liberté''
sur le versant yang, en faisant mine d'ignorer superbement ``le chaos''.}

\deleted{Cela m'a fait constater tout d'abord que parmi les quarante qualités
et entités différentes qui figurent soit dans la Fleur cosmique, soit
dans l'un ou l'autre de nos deux ``paquets attracteurs'', il y en a
d\,e\,u\,x exactement qui se distinguent de tous les autres comme des sortes
d' ``a\,n\,t\,i\,-\,a\,t\,t\,r\,a\,c\,t\,e\,u\,r\,s'' , pour ne pas dire, de
``repoussoirs''. Ce sont justement deux parmi les trois épouses de l'ordre
(résolument polygâme...) dont nous avons fait connaissance jusqu'à présent\handwritten{,}
savoir les dames}

\[
\text{\deleted{chaos , \quad liberté}}
\]

\deleted{-- à l'exclusion donc de \deleted{``mystère'', qui} dame mystère, qui elle, par
contre, \deleted{\unsure{représente}} exerce une attraction de force exceptionnelle.}

\deleted{Il semblerait que liidée même du chaos suscite dans l'esprit humain
une répugnance quasi\deleted{\unsure{-}} insurmontable. Nous le ressentons, ``viscéralement'',
comme chose irréductiblement o\,p\,p\,o\,s\,é\,e à l'ordre, objet de}

\scanpage{130}
% archive sheet number: 129

Supposons maintenant que \deleted{$H,H'$} l'on
se donne une alg. Azumaya $A$ et des
plongements

\margin{plus \unsure{gén.|génl}, dans une $A$ admettant \unsure{formules|formule|fournissant} H.C. d'où $\Delta_A,\mathrm{Tr}_A$ .....}

\diagram{\begin{tikzcd}[ampersand replacement=\&,row sep=small]
\& A \& \\
H \arrow[ur,"i"] \& \& H' \arrow[ul,"j"']
\end{tikzcd}}

On peut alors définir hom de $k$-modules

\diagram{\begin{tikzcd}[ampersand replacement=\&]
P \arrow[r,"\deleted{\unsure{t}}f"] \& A
\end{tikzcd}}

par la condition

\[
 f(\deleted{\unsure{x}}u\otimes v)=i(u).j(v)
\]

\deleted{ce qui induit}
et on récupère
\deleted{et la donnée de}
$i,j$ par

\[
 i=f\circ i_0,\quad j=f\circ j_0
\]

La donnée de $i,j$ équivaut à la
donnée de $f$, soumis aux
conditions que $f\mid H$ et $f\mid H'$
soient des hom d'algèbres. On a alors

\diagram{\begin{tikzcd}[ampersand replacement=\&,row sep=small,column sep=small]
\& A \& \\
\& P \arrow[u,"\unsure{f|t}"'] \arrow[d] \& \\
\& B \arrow[dl] \arrow[dr] \& \\
H \arrow[uur, crossing over, "\deleted{\nearrow}" description,near end] \& \& H' \\
\& k \arrow[ul,"\varepsilon_H"] \arrow[ur,dash,"\varepsilon_{H'}"'] \&
\end{tikzcd}}

\margin{attention, on ne suppose pas que $f$ soit un iso de Mod i.e que \deleted{\unsure{$A$}}$H,H'$ engendrent $A$ comme algèbre}

\scanpage{131}
% archive sheet number: 130

Cette situation permet de définir
sur $P$

\[
\begin{aligned}
&\text{forme linéaire }L_P=\mathrm{Tr}_A\circ f,\ \text{qui prolonge }T_B\\
&\text{forme quadratique }\Delta_P=\det_A\circ f\\
&\text{forme bilinéaire sym. }\psi_P\ \text{polarisée de }\Delta_P,\ \psi_P=\psi_A\circ f\\
&\hspace{6em}\text{(où }\psi_A\text{ est polarisée de }\det_A\text{)}\\
&\text{forme bilinéaire }\varphi_P=\varphi_A\circ f
\end{aligned}
\]

i.e.

\[
\varphi_P\deleted{\unsure{=}}(U,V)=\mathrm{Tr}\bigl(f(U)\unsure{,|.}f(V)\bigr)
\]

\margin{\uline{NB} On a \deleted{$\unsure{\Delta_P,\varphi_P,\psi_P}$} $\varphi_P(U,V)+\psi_P(U,V)=L_P(u)L_P(V)$ et $\varphi_P(1,U)=\psi_P(1,U)=L_P(U)$.}

D'autre part, on peut restreindre ceci
à $B$, on trouve

\[
\text{forme quad }\Delta_B,\quad\text{formes bil. sym }\psi_B,\varphi_B.
\]

En restreignant ces trois formes sur
$B$ à $H$, on trouve les trois formes
dont $H$ est déjà équipé, $\Delta_H=\det_H$, $\psi_H=$ forme
polarisée de $\Delta_H$, $\varphi_H=\bigl((u,v)\unsure{-|\mapsto}\mathrm{Tr}_H(uv)\bigr)$, Et
itou pour $H'$.

On a d'autre part

\[
\Delta_P(u\otimes v)=\Delta_H(u)\,\Delta_{H'}(v) \tag*{(1)}
\]

\uline{Proposition} \inserted{On se donne $H,H'$ comme dessus.} Considérons les \deleted{\unsure{trois|six}} \inserted{quatre} ensembles
suivants :

$Q_P$: \deleted{$(\unsure{\Delta_P})$} ens des formes quadratiques $\Delta_P$ sur $P$ satisfaisant
(1) ci-dessus

$Q_B$: \deleted{$(\unsure{\Delta_B|\Delta_H})$} ens des formes \deleted{\unsure{bil.}} \unsure{qdr.|quadr.} $\Delta_B$ sur $B$ qui
prolongent $\Delta_H$ et $\Delta_{H'}$

$\underline{\Psi}_B$: \deleted{$(\unsure{\psi})$} ens des formes bil. sym $\psi_B/\Delta_B$ sur $B$
qui prolongent $\psi_H,\psi_{H'}$

$\underline{\Phi}_B$: \deleted{$(\varphi)$} ens des formes bil. symétriques $\varphi_B$ sur $B$

\scanpage{132}
% archive sheet number: 131

2) \deleted{On voit que} \inserted{Posons}

\[
L=H/k\inserted{\simeq B/H'},\quad L'=H'/k\simeq B/H
\]

\deleted{et considérons une forme}

\[
\deleted{\lambda:L\otimes L'\unsure{-|\longrightarrow} k}
\]

\deleted{Si $\Delta_B\in\underline{Q}_B$, on peut définir}

On a

\[
\begin{aligned}
\deleted{\unsure{\Delta_B\to B}}\ \deleted{\unsure{B\to L\otimes{}}}\quad
(L\otimes L')^\vee&\xrightarrow{\sim}\text{Esp. des formes}\\
&\quad\text{bil. \inserted{sym} sur }B\text{ nulles}\\
&\quad\text{sur }H'\deleted{\unsure{B}}\text{ et }H\deleted{\unsure{B}}
\end{aligned}
\]

(i.e. qui passent au
quotient dans $B/k\simeq L\oplus L'$ \inserted{car $B=H+H'$}
et qui sont nulles
sur $L$, et sur $L'$
donc de la forme

\[
\begin{aligned}
(\underset{\rotatebox{90}{$\in$}\,L}{x},\underset{\rotatebox{90}{$\in$}\,L'}{y}\deleted{\unsure{'}}),\
\deleted{(\unsure{x',y'})}\inserted{(\xi,\unsure{\eta|y'})}
 &\longrightarrow\lambda(x\otimes\unsure{\eta|y})\\
 &\quad+\lambda(\xi\otimes y)
\end{aligned}
\]

avec $\lambda:L\otimes L'\longrightarrow k$

Donc \deleted{$\Delta$}$\unsure{\underline{\Phi}_B|\Phi_B}$ est un torseur sous $(L\otimes L')^\vee$, et de
même pour $\underline{\Psi}_B$. Itou pour $\Delta_B$, car

\[
\begin{aligned}
(L\otimes L')^\vee&\simeq\text{formes quadratiques sur }B\\
&\qquad\text{qui sont nulles sur }H,H'
\end{aligned}
\]

(donc passent au quotient
dans $B/k\simeq L\oplus L'$ en des
formes quadratiques nulles
sur $L,L'$ \deleted{\unsure{donc formes de|la forme}} donc de
la forme

\[
\Delta(x\oplus y)=\lambda(x\otimes y)
\]

Les applications $\beta,\gamma$ sont compatibles
avec ces structures de torseurs. Cela
suffit à assurer qu'elles sont iso.

\scanpage{133}
% archive sheet number: 132

qui prolongent $\varphi_H,\varphi_{H'}$.

$\deleted{\unsure{T}}L_P$: Ens des formes \uline{lin.} sur $P$ qui prolongent $T_B$

On a des flèches canoniques

\diagram{\begin{tikzcd}[ampersand replacement=\&,row sep=tiny,column sep=small]
\deleted{\unsure{Q_P|\underline{\Phi}_P}}\ \deleted{\underline{\Phi}_P}\arrow[r,"{\deleted{\to}}" description] \& \deleted{\underline{\Phi}_B}\arrow[r,"{\deleted{\to}}" description] \& \deleted{\unsure{\underline{\Phi}|\underline{\Phi}_B}}
\end{tikzcd}}

\diagram{\begin{tikzcd}[ampersand replacement=\&,row sep=small,column sep=small]
\underline{Q}_P\arrow[r,"\alpha"]\arrow[dr,"\delta",swap]\arrow[dr,phantom,shift right=2ex,"{\deleted{\unsure{\varphi}}}"'] \& \underline{Q}_B\arrow[r,"\beta"] \& {\deleted{\unsure{\underline{\Phi}_B}}\,\underline{\Psi}_B}\arrow[r,"\gamma"] \& \underline{\Phi}_B\arrow[dll,"\zeta"]\\
\& \underline L_P \& \deleted{\unsure{\varphi_P\to\underline{\Phi}_B}} \&
\end{tikzcd}}

que voici : $\alpha$ est la restriction $\alpha(\deleted{\unsure{q}}\Delta_P)=\deleted{\unsure{q}}\Delta\circ\alpha$,
$\beta$ associe à tte forme $\Delta_B\in Q_B$ la
forme bil. sym. associée \inserted{$\psi_B$}, enfin
on définit $\gamma$ par la formule

\[
\underset{\varphi_B}{\underbrace{\gamma(\psi)}_{\parallel\,\mathrm{déf}}}
\deleted{\unsure{=}}+\psi_B=\deleted{\unsure{\{}}
\bigl((u,v)\longmapsto T_B(u)T_B(v)\bigr)
\]

Enfin, on définit $\delta$ par

\[
\delta(\Delta_P)=\{u\longmapsto\varphi_{\Delta_P}(u,1_P)\}
\]

% slanted
\margin{On définit $\zeta:\underline{\Phi}_B\to\underline L_P$ ainsi : $\underline L_P$ s'identifie \inserted{à l'}\deleted{\unsure{au|avec|aux}}ens des formes \uline{bilin.} sur $H\times H'$ telles que leurs restrictions à $1\times H'$, et à $H\times1$ sont données par $T_{H'}$ et $T_H$. à $\varphi_B$ on associe la forme bil. $\zeta(\varphi_B)\in\unsure{*|:}\,(u,v)\longmapsto\varphi_B(u,v)$. C'est encore un hom de torseurs sous $(L\otimes L')^\vee$, donc iso.}

Ceci dit, les quatre applications
$\alpha,\beta,\gamma,\delta$ sont \uline{bijectives}.

\uline{Remarques} 1) Pour définir les cinq
ensembles $\underline Q_P,\underline L_P,\underline Q_B,\underline\Psi_B,\underline\Psi_B$, \deleted{et}
les applications $\alpha,\beta,\gamma,\delta$, on ne
utilise sur $H,H'$ que la donnée
de $\Delta_H,\Delta_{H'}$ et $1_H,1_{H'}$ ($\unsure{\longrightarrow|\Longrightarrow}T_H,T_{H'}$ \unsure{etc|cf.})
satisfaisant $\Delta_H(1_H)=\Delta_{H'}(1_{H'})=1$.

\scanpage{134}
% archive sheet number: 133

$\simeq(\det H)\otimes\det(H')$. Cette \deleted{\unsure{loi}} fonction
quadratique est \uline{universelle}, pour
les fonctions quadratiques

\diagram{\begin{tikzcd}[ampersand replacement=\&]
P\arrow[r,"q"] \& Q
\end{tikzcd}}

telles que $q(x\otimes y)=0\qquad\forall x,y\in P$.

Ceci nous montre \deleted{que, si} (revenant
à \deleted{$\unsure{Q_P|\underline Q_P}$} la situation \unsure{précédente|précédent}) que
si $\underline Q_P$ n'est pas vide, c'est un torseur
sous \deleted{$\unsure{\underline L_B|\underline L_P}$} $(L\otimes L')^\vee$. Il faut vérifier
que $\alpha,\delta$ sont compatibles avec
structure torseur -- moyennant quoi
on voit que c'est des iso.

Pour vérifier que $\underline Q_P$ n'est pas
vide, on \deleted{prend sur $P$ la structure
\deleted{\unsure{tensorielle}} d'alg. produit tensoriel
de celles de $H,H'$, \deleted{et la fon} trouve une Alg. de H.C., dont
la fonction $\det$ fait l'affaire.}

on peut utiliser la construction
des algèbres $M(b,c,b',c',\underline t)$. Cela
me fait songer à un sixième
ensemble

\scanpage{135}
% archive sheet number: 134

en fait, comme on a

\[
H\unsure{\longrightarrow|-}L,\quad H'\unsure{\longrightarrow|-} L'
\]

on en déduit

\[
P\overset{\mathrm{dfn}}{=}H\otimes H'\longrightarrow L\otimes L'\qquad\text{épi}
\]

nul sur $1\otimes H'$ et sur $H\otimes1$, donc sur
$B$, et on a la suite exacte

\diagram{\begin{tikzcd}[ampersand replacement=\&,column sep=small]
0\arrow[r] \& B\arrow[r] \& P\arrow[r] \& \deleted{\unsure{1}}L\otimes L'\arrow[r] \& 0
\end{tikzcd}}

\deleted{Ceci dit}
Ainsi les formes linéaires \inserted{$T_P$} sur
$P$ dont la restr. à $B$ est $T_B$, font
un torseur \inserted{$\underline L_P$} sous $(L\otimes L')^\vee$, L'application
can.

% slanted
\margin{i.e. les formes bil. $H\times H'\xrightarrow{f}k$, telles que $f(u,1)=\mathrm{Tr}_H(u)$, $f(1,v)=\mathrm{Tr}_{H'}(v)$}

\[
\begin{aligned}
\zeta:\ \underline{\Phi}_B&\longrightarrow\underline L_P\\
\varphi_B&\longmapsto (x\otimes y\longmapsto\varphi_B(x\deleted{\unsure{\otimes1,1\otimes{}}},y))
\end{aligned}
\]

est compatible avec les opérations
de $(L\otimes L')^\vee$. Donc $\zeta$ est iso.

Notons le

\uline{Lemme} Soient $H,H'$ deux modules vectoriels
de rang 2 sur $k$, $P\simeq H\otimes H'$. On définit
une fonction quadratique canonique

\diagram{\begin{tikzcd}[ampersand replacement=\&]
P\arrow[r,"{\unsure{\delta_{H,H'}|q_{H,H'}|\Delta_{H,H'}}}"] \& \underbrace{\det H}_{L}\otimes\underbrace{\det H'}_{L'}
\end{tikzcd}}

en interprétant $P\simeq\mathrm{Hom}(H^\vee,H')\longrightarrow\mathrm{Hom}(\det H^\vee,\det\unsure{H'})$

\scanpage{136}
% archive sheet number: 135

$\mathrm{Az}_P=$ structure de $k$-algèbre de H.C.
sur $P=H\otimes H'$, induisant sur $1\otimes H'$
et $H\otimes1$ les structures d'alg. données.

\diagram{\resizebox{\linewidth}{!}{\begin{tikzcd}[ampersand replacement=\&,column sep=large,row sep=huge]
\mathrm{Az}_P\arrow[r,"{\substack{\alpha\mapsto\det_\alpha\\\deleted{\unsure{\alpha\mapsto\det_\alpha}}}}"]
\arrow[drr,"{\alpha\longrightarrow\mathrm{Tr}_\alpha}"',bend right=7]
\& Q_P\arrow[r,"{\Delta_P\mapsto\Delta_P\mid B}"]
\arrow[dr,dashed,"{\Delta_P\mapsto(U\mapsto\psi_{\Delta_P}(1,U))}"]
\arrow[dr,phantom,"{\text{\unsure{comm.}}}"',near start]
\& Q_{\underline B}\arrow[r,"{\Delta_B\to\varphi_{\Delta_B}}"]
\& \Psi_{\underline B}\arrow[r,"{\varphi\mapsto\psi\ \text{t.q.}\ \varphi+\psi=T_B\otimes T_B}"]
\& \Phi_B\arrow[dll,"{\varphi\mapsto(\underset{\substack{\rotatebox{90}{$\in$}\\H}}{x}\otimes\underset{\substack{\rotatebox{90}{$\in$}\\H'}}{y}\mapsto\varphi(x,y\unsure{'|{}}))}"']\\
\& \& L_P \& \&
\end{tikzcd}}}

La commutativité du triangle total
est immédiate, ainsi que celle du
triangle de gauche, on en conclut
celle du triangle de droite (qui
n'était pas évidente sinon).

\deleted{L'on} On aurait pu aussi
introduire $\underline\Psi_P$ et $\underline\Phi_P$ (en les définissant
avec soin) mais passons ! Parmi les
6 \deleted{\unsure{bases}} ens, celui qui est de
description alg. la plus simple est
sans doute $L_P\subset\mathrm{Bil}(H,H';k)$. On
appellera un $T_P\in L_P$ un \uline{twist}\deleted{\unsure{ing}}
des algèbres \deleted{\unsure{d'Az}} quadratiques
$H,H'$.

% slanted
\margin{C'est défini en termes de $H,H'$ munis \unsure{seulemt|seulement} de $1_H,1_{H'}$ et de $T_H,T_{H'}$, à l'exclusion de la donnée de $\Delta_H,\Delta_{H'}$}

L'objet le \uline{plus fort}
qui en résulte, est la structure
d'alg. de H.C. sur $P=H\otimes H'$,

\scanpage{137}
% archive sheet number: 136

qui donne aussitôt les éléments
de structure

\[
\Delta_P,T_P,\varphi_P,\psi_P,\deleted{\unsure{\varphi}}\Delta_B,\psi_B,\varphi_B
\]

\uline{Th. récapitulatif} \textcircled{1} Les flèches
du diagramme de la page précédente
sont bijectives, le diagramme est
comm, et les \inserted{5} flèches qui ne sont
pas de source $\mathrm{Az}_P$ sont \deleted{\unsure{don}}compatibles
avec les structures de torseur
sous $(L\otimes L')^\vee$.

\textcircled{2} Pour qu'une structure \inserted{d'alg sur $P$} associée
à un twist en fasse une Alg.
d'Azumaya, il f et s. que
la forme $\Delta_B$ sur $B$ soit non
dégénérée, i.e que son discriminant
divisé

\[
\mathrm{disc}'(\Delta_B)\in(\det B)^{\otimes-2}=(L\otimes L')^{\otimes-2},
\]

soit inversible.

% slanted
\margin{Source \deleted{\unsure{et|cette}} et but ne dépendent que de $1_H,1_{H'},T_H,T_{H'}$ mais l'application $\delta$ \uline{fait} appel à $\Delta_H,\Delta_{H'}$}

\textcircled{3}

\diagram{\begin{tikzcd}[ampersand replacement=\&,column sep=huge]
\underline L_P\arrow[r,"{\delta_{P,B}}","{\substack{\text{discriminant}\\\text{divisé de}\\\text{la forme quad. sur}\\ B\ \text{associée}\\\text{à un twist}}}"']
\& (L\otimes L')^{\otimes-2}
\end{tikzcd}}

est une application \deleted{quadratique} pol. de degré $2$

\scanpage{138}
% archive sheet number: 137

dont la forme quadratique associée

\diagram{\begin{tikzcd}[ampersand replacement=\&,column sep=large]
(\underbrace{L\otimes L'}_{\substack{\text{esp. des}\\\text{translations}\\\text{de }\underline L_P}})^{-1}
\arrow[r,"{\deleted{\unsure{\delta_{P,B}^{0}|\delta_{P,B}}}}"] \& (L\otimes L')^{-2}
\end{tikzcd}}

est l'élévation au carré (donc $\delta_{P,B}$
est \deleted{qui pour la} un morphisme fini
surjectif de degré 2).

Prenons l'image inverse de la section
nulle, c'est un schéma quadratique
\deleted{\unsure{$\widetilde S$|fini}} sur $S=\mathrm{Spec}\,k$ (donc défini par Alg.
quadratique $k'$) \inserted{qui ne dépend que des \unsure{rev.|inv.|env.} quad.}, cela dit, on a
un \uline{morphisme canonique}

% slanted
\margin{$X,X'$ \unsure{de|des} $S$ définis par $H,H'$, \unsure{soient|soit|posé} $\widetilde S=X\wedge X'$}

\diagram{\fbox{\begin{tikzcd}[ampersand replacement=\&]
X\times_S X'\arrow[r,"\alpha"] \& \widetilde S
\end{tikzcd}}}

\deleted{\unsure{ce qui|tel que}} tel que

\[
\alpha(\bar x,x')=\alpha(x,\overline{x'})=\overline{\alpha(x,x')}
\]

(où $x\unsure{-|\mapsto}\bar x$ etc désigne les inv. canoniques)
d'où pour $X,X'$ étales (et alors $\widetilde S$ est étale)

\diagram{\begin{tikzcd}[ampersand replacement=\&]
\overbrace{X\wedge_S X'}^{\substack{\text{produit contracté}\\\text{de }(\mathbb Z/2\mathbb Z)_S\text{-torseurs}}}
\arrow[r] \& \widetilde S,
\end{tikzcd}}

\deleted{\unsure{défini}} hom de $(\mathbb Z/2\mathbb Z)_S$-torseurs, donc \uline{iso}.

Il faut définir l'application $\alpha$

Soient $\varphi:\deleted{\unsure{A}}H\unsure{-|\to} k,\quad \varphi':H'-k$

\deleted{deux hom. d'Alg, considérons la forme}

\[
\deleted{\varphi\otimes\varphi':H\otimes H'=P\longrightarrow k,\quad u\otimes v\mapsto\unsure{\varphi(u),\varphi(v)|\varphi(u)\varphi(v)|\varphi(u)\varphi'(v)}} 
\]

\scanpage{139}
% archive sheet number: 138

Choisissons bases $\unsure{1(,u)|(1,u)}$ de $H$, $(1,v)$ de $H'$
(d'où $b,c,b',c'$), $\varphi$ et $\varphi'$ sont données par

\[
\varphi(u)=z,\qquad\varphi'(v)=z'
\]

avec $z\in k$

\[
z^2-bz+c=0\qquad z'{}^2-b'z+c'=0,
\]

et on cherche $t$ solution de

\[
t^2-bb't+(b^2b'+b'{}^2b-4cc')=0.
\]

Mais le discriminant de cette équation est

\[
\Delta=(\deleted{\unsure{b}}bb')^2-4\bigl(b^2\deleted{b}c'+b'{}^2\deleted{b}c-4cc'\bigr)=
\underbrace{(b^2-4c)}_{\delta}\underbrace{(b'{}^2-4c')}_{\delta'}
\]

Or \deleted{\unsure{comme}} (\unsure{se|en|les} plaçant \deleted{en car. 2} heuristiquement
d'abord sur un corps de car $\ne2$)

\[
2z=b\pm\sqrt{\delta},\qquad 2z'=b'\pm\sqrt{\delta'}
\]

\[
\text{\deleted{donc}}\qquad\delta=(2z-b)^2\qquad\delta'=(2z'-b)^2
\]

(vrai sans condition corps ou caract.)
et les solutions de l'équation en $t$ sont \inserted{données} par
(sur corps car $\ne2$) \unsure{d}

\[
\deleted{\unsure{\tfrac14}}\,(2t-bb')^2=\Delta=\delta\delta'=\bigl((2z-b)(2z'-b')\bigr)^{\unsure{2}}
\]

donc on peut \deleted{poser} supposer

\[
2t-bb'=\pm(2z-b)(2z'-b')
\]

soit

\[
2t=bb'\pm(2z-b)(2z'-b').
\]

Choisissant $-$ des \deleted{\unsure{sig}} les $\pm$, on trouve

\[
2t=\deleted{\unsure{4}}-4zz'+2(\deleted{\unsure{4}}b'z+bz')\deleted{\unsure{+}}\qquad\text{\deleted{soit}}
\]

\[
\boxed{t=(b'z+bz')-2zz'}
\]

l'autre racine sera

\[
t^{\unsure{1|'}}=\bigl(bb'-(b'z+bz')\bigr)+2zz'\ (=bb'-\deleted{\unsure{z}}t)
\]

\scanpage{140}
% archive sheet number: 139

Il suffit de vérifier que l'on a bien

\[
tt'=b^2c'+b'{}^2c-4cc'\qquad\text{(\unsure{vu|on} qu'on a déjà }t+t'=bb'\text{)}
\]

Vérification immédiate (il n'est pas besoin de \unsure{le|la} faire)
On peut l'écrire

\[
\boxed{\deleted{\unsure{t}}T_{\varphi,\varphi'}\deleted{\unsure{(u,v)}}\unsure{(u,v)}=\mathrm{Tr}(v)\varphi(u)+\mathrm{Tr}(u)\varphi'(v)-2\deleted{\unsure{T}}\varphi(u)\varphi'(v).}
\]

Or, on a là une forme bil. en $u,v$,
qui pour $v=1$ se réduit :

\[
t(u,1)=\deleted{\unsure{\mathrm{Tr}}}2\varphi(u)+\mathrm{Tr}(u)-2(\varphi)=\mathrm{Tr}(u)
\]

et de \unsure{$=$|même} pour $u=1$ \deleted{\unsure{on}}

\[
t(u,1)=\mathrm{Tr}\,u,\qquad t(1,v)=\mathrm{Tr}\,v
\]

C'est donc un élément de $L_P$ ! Soit
\deleted{Je dis que $\unsure{\delta}$}
$\deleted{\unsure{t}}T_{\varphi,\varphi'}\in L_P$
défini par $\varphi,\varphi'$. J dis que

\[
\delta(t_{\varphi,\varphi'})=0
\]

Pour le voir, on peut faire le calcul
avec bases \deleted{\unsure{quelc.|quelconque}} $(1,u)$ $(1,v)$ telles que

\[
\varphi(u)=\varphi(v)=0.
\]

\inserted{d'où $\overset{c}{\overset{\parallel}{\det}}u=\overset{c'}{\overset{\parallel}{\det}}v=0$ car \deleted{\unsure{on prend}} $\varphi(u)=0$ \unsure{sol. de|de} $z^2-bz+c=0$}

de sorte que maintenant

\[
\deleted{\unsure{b'}}t(u,v)=0\deleted{\unsure{2\varphi}}.
\]

\deleted{Je dis que l'équation} Or on a l'identité (

\[
\begin{aligned}
\delta(u,v,t)&=t^2-bb't+(b^2c'+b'{}^2c-4cc')=\deleted{\unsure{t^2-bb't}}\\
&=t(t-bb')\qquad\text{car }c=c'=0\ \deleted{\text{car }c=c'=0}
\end{aligned}
\]

donc \deleted{en y faisant $t$} $t(u,v)=0$ en est
bien une solution.

\scanpage{141}
% archive sheet number: 140

\deleted{Notons} On a bien défini ainsi

\diagram{\begin{tikzcd}[ampersand replacement=\&]
X\times X'\arrow[r] \& X\overset{\deleted{\unsure{S}}}{\wedge}X'\quad(\overset{\mathrm{déf}}{=}\widetilde S)
\end{tikzcd}}

(car la définition précédente commute
au chg\textsuperscript{t}. de base).

Notons qu'on a dans $L_{\underline P}$ une involution
canonique,

\[
\deleted{\unsure{\varphi}}\,T\longmapsto T'
\]

où

\[
T'(u,v)=T(u',v)=T(u,v')
\]

(Il faut prouver que l'on a bien

\[
\boxed{T(u',v)=T(u,v')}\ \overset{\mathrm{déf}}{=}T'(u,v)
\]

pour tte forme $T\in\underline L_P$, \deleted{\unsure{cela|alors en}} en
termes des plong\textsuperscript{ts} $H,H'\subset A_T$,
cela signifie simplement

\[
\mathrm{Tr}_{A_T}(u'v)=\mathrm{Tr}_{A_T}(\underbrace{uv'}_{\deleted{\unsure{(uv)'}}})
\]

or le premier mb, compte tenu que
$\mathrm{Tr}_{A_T}$ est \unsure{invariant|invariante} sous $'$, est aussi
$\mathrm{Tr}_{A_T}\bigl((u'v)'=v'u\bigr)$ et $\mathrm{Tr}\,v'u=\mathrm{Tr}\,uv'$ OK.

On a donc

\[
(T')'=\deleted{\unsure{\varphi}}T,
\]

et on voit aussitôt que

\scanpage{142}
% archive sheet number: 141

\[
\delta(T)=\delta(T')
\]

donc $T\mapsto T'$ induit un \deleted{\unsure{$k$}} autom
de $\widetilde S=X\wedge X'=\delta^{-1}(\unsure{-|{}}\text{section nulle de }\unsure{\check V|V^\vee|\mathbb V}(L\otimes L'))$.

Quand on choisit bases $u,v$ donc
$L_P$ identifié à l'axe des $t$, l'involution
$T\mapsto T'$ \deleted{à qui} se représente par

\[
t\longmapsto bb'-t,
\]

qui \deleted{\unsure{est|$\Rightarrow$}} \inserted{induit} l'\uline{involution can.} sur $\widetilde S$, schéma
des solutions de

\[
t^2-bb't+\text{fourbis}=0.
\]

\deleted{\unsure{A|$\triangle$}} On va donner une interprétation
schématique de $X\wedge X'$ comme produit
contracté.

\uline{Prop} \deleted{\unsure{Soit}}Supposons que 2 soit inv., ou
plus généralement, qu'en tt pt où 2 est nul,
soit $X$ soit $X'$ est étale (cela signifie
aussi que $2,\underset{\substack{\parallel\\ b^2-4c}}{\deleted{\unsure{b}}\delta},\underset{\substack{\parallel\\ b'{}^2-4c'}}{\deleted{\unsure{b'}}\delta'}$ ne s'annulent pas
simultanément sur $S$, i.e que
l'idéal $(2,\delta,\delta')$ de $k$ est l'idéal unité).
Alors le \deleted{\unsure{hom}} morphisme

\diagram{\begin{tikzcd}[ampersand replacement=\&]
X\times_S X'\arrow[r] \& X\wedge_S X'
\end{tikzcd}}

est \deleted{\unsure{içi fid|fid}} un morphisme quadratique
(fini loc\textsuperscript{t} libre de rang 2), et si $\sigma_X,\sigma_{X'}$ sont
les inv can. de $X$ et $X'$ resp\textsuperscript{t}, $\sigma_X\times_S\sigma_{X'}$ est l'inv. can

\scanpage{143}
% archive sheet number: 142

de $X\times_S X'$ sur $X\wedge_S X'$, et celui-ci est
le \deleted{sous} schéma quotient (universel)

\[
X\wedge_S X'\simeq X\times_S X'\big/\underbrace{\sigma_X\times_S\sigma_{X'}}_{=\sigma_P}
\]

conformément au yoga général des
produits contractés. Pour que
$X\times_S X'\unsure{-|\longrightarrow}X\wedge_S X'$ soit étale, il f. et s. qu'en
tt $s\in S$, $X_s$ \uline{ou} $X'_s$ soit étale. Pour
que $X\wedge_S X'$ soit étale sur $S$, il f et s.
que \deleted{en tout point $s\in\unsure{S|k}$}, $X$ \inserted{\uline{et}} $X'$ soient étales.

(Dans ce cas, on retrouve la notion familière.
Si $X$ \unsure{seult|seulemt|seule} est étale, on peut considérer
$X\wedge X'$ comme ``twisté'' de $X'$ par le
$(\mathbb Z/2\mathbb Z)_S$-torseur $X$ et l'op. de $\mathbb Z/2\mathbb Z$ sur $X'$
par son inv. can.)

En termes des Alg. $H,H'$ de $X,\unsure{x'|X'}$,
et $P=H\otimes H'$ de $X\times_S X'$ (avec cette fois
la structure d'Alg. comm. can. sur $P$ !]
on peut décrire celle de $X\wedge X'$ comme
la sous-Alg. \inserted{$H\wedge H'$} des invariants sous $'$.
Si $u,v$ sont des bases, alors $H\wedge H'$
a comme base $1\ (=1\otimes 1)$ et

\[
\boxed{u*v\overset{\mathrm{déf}}{=}\mathrm{Tr}(v)u+\mathrm{Tr}(u)v-2\underbrace{uv}_{\substack{\parallel\\ u\otimes v}}}
\ \in(H\otimes H')^{\unsure{\sigma_P}}
\]

% slanted
\margin{\uline{NB} Cet élément est défini sans \unsure{aucunes|aucune} conditions (\deleted{\unsure{de}} car etc) sur $H,H'$ \deleted{\unsure{ou}} \unsure{ou} $u,v$}

On a alors

\scanpage{144}
% archive sheet number: 143

\[
\left\{\begin{aligned}
\mathrm{Tr}^{\unsure{q|\circ|-}}(u*v)&=bb'\\
\det(u*v)&=\delta(u,v)=b^2c+b^2c'-4cc',
\end{aligned}\right.
\]
\margin{il s'agit ici de la trace et du déterminant en tant qu'élément d'une Alg. quadratique $H\wedge H'$ sur $S$}
% slanted
\margin{\uline{NB} Si $u,v$ sont des \uline{bases}, alors $u*v$ n'est autre que la \unsure{restriction} à $\widetilde S\subset\check V(L_P)$ de la fonction paramétrique $t_{u,v}$ sur $\check V(L_P)$ \unsure{Aff.|Ainsi,} Si on considère $u*v$ comme \deleted{\unsure{la}} \unsure{fct}, \deleted{\illegible{1}} \deleted{\unsure{bil.}} \deleted{\illegible{1}} sur $H^\vee\otimes H'^\vee$, sa restriction à $L_P\subset H^\vee\otimes H'^\vee$ est la fonction paramétrique $t_{uv}$ sur $L_P$.}

Ainsi, la formule \deleted{donnée} plus haut
\deleted{\unsure{donn.}} pour $T_{\varphi,\varphi'}(u,v)$ peut s'écrire

\[
\boxed{T_{\varphi,\varphi'}(u,v)=\varphi\otimes\varphi'\,(u*v)}\,;
\]

Notons aussi les formules

\[
\underset{\text{discr}}{\delta}(u*v)=\delta(u)\delta(v)
\]

ainsi que

\[
\left\{\begin{aligned}
\underbrace{(u*v)'}_{\substack{\text{le $'$ est relatif}\\\text{à l'alg. quad}\\ H\wedge H'\,\deleted{\unsure{,}}}}&=u'*v=u*v'\\
u'*v'&=u*v
\end{aligned}\right.
\]

\uline{NB} Toutes les formules précédentes ont un
sens et sont valables, sans restriction
de car. sur $S,H,H'$, pour \uline{tt} $u\in H$ et $v\in H'$
(base ou pas !)

On récupère $L_{\deleted{\unsure{\varphi}}P}\overset{\text{dfn}}=L_{H,H'}$ à partir de
$\widetilde S=X\wedge X'$ comme l'espace affine enveloppant
(au dessus de $S$) du rev. quad. $\widetilde S$, et
$\delta:L_P\longrightarrow(L\otimes L')^{-1}$ comme l'hom (discriminant !) canonique associé à \deleted{\unsure{le|ce}} revêtement.

